\documentclass[10pt]{article}

\usepackage[utf8]{inputenc}

\usepackage[T1]{fontenc}

\usepackage{amsmath}

\usepackage{amsfonts}

\usepackage{amssymb}

\usepackage[version=4]{mhchem}

\usepackage{stmaryrd}

\usepackage{bbold}

\usepackage{graphicx}

\usepackage[export]{adjustbox}

\graphicspath{ {./images/} }



\begin{document}

$\langle\boldsymbol{M}\rangle_{V}=\boldsymbol{e} M_{v} \neq 0$ при $v \neq 0$. Квазисреднее магнитного момента равно



\[

\lim _{v \rightarrow 0} M_{v}

\]



и отлично от нуля при температурах ниже точки Кюри. При построении К. существенно, что $v \rightarrow 0$ после выполнения термодинамич. предельного перехода $V \rightarrow \infty$ при фиксированном $V / N$, где $N-$ число частиц. Если эти предельные переходы перестановочны, то К. равны нулю, это справедливо при температурах выше точки Кюри. В общем случае квазисреднее оператора $A$ равно



\[

\lim _{v \rightarrow 0}\langle A\rangle_{v}

\]



где $\langle A\rangle_{v}-$ обычные статистич. средние при наличии поля ve, снимающего вырождение. Обычные средние равны К., усредненным по всем направлениям поля. Аналогично вводят К. в теории кристаллов, нарушая симметрию, связанную с пространственными трансляциями и вращениями, в теории сверхтекучести и сверхпроводимости, где нарушают симметрию гамильтониана, связанную с сохранением полного числа частиц; в квантовой теории поля и т. д.



Общий способ введения К. таков. Рассматривают макроскопич. систему с гамильтонианом $H$. Добавляют к $H$ бесконечно малые добавки, нарушающие нек-рые законы сохранения (симметрию гамильтониана), получая гамильтониан $H_{v}$. Если все средние значения $\langle A\rangle_{v}$ получают лишь бесконечно малые приращения, состояние статистич. равновесия называется не вырожденным. Если же нек-рые из средних получают конечные приращения, говорят о вы рожде н и и состояния статистич. равновесия. В этом случае вводят К., равные



\[

\lim _{v \rightarrow 0}\langle A\rangle_{v}

\]



причем сначала выполняется предельный переход $V \rightarrow \infty$.



К. удобны также для вычисления корреляционных функций, функций Грина и т. П.



Обычный метод теории возмущений, строго говоря, не применим к системам с вырожденным состоянием статистич. равновесия. Для того чтобы воспользоваться теорией возмущений в этом случае, нужно предварительно снять вырождение и ввести функции Грина, построенные из К.



Лum.: [1] Ахи е зе р А. И., П е ле т м и н ск и й С. В., Методы статистической физики, М., 1977; [2] Б о го л ю бо в Н. Н., Избранные труды по статистической физике, М., 1979; [3] Б о го л ю бо в Н. Н., Бо гол юб о в Н.Н. (мл.), Введение в квантовую статистическую механику, М., 1984.



Д. Н. Зубарев.



КВАЗИЧАСТИЧЫ - элементарные возбуждения конденсированной среды (твердого тела, жидкого гелия), ведущие себя в некоторых отношениях как квантовые частицы. Теоретич. описание и объяснение свойств конденсированных сред, исходящее из свойств составляющих их молекул, атомов, ионов и электронов, представляет большие трудности, во-первых, потому что число частиц огромно $\left(10^{22}-10^{23}\right.$ в $1 \mathrm{~cm}^{3}$ ), а, во-вторых, потому что частицы сильно взаимодействуют между собой. Из-за взаимодействия частиц полная энергия конденсированной среды не есть сумма энергий отдельных частиц (как в идеальном газе). Развитие квантовой теории твердого тела привело к концепции К., к-рая оказалась особенно плодотворной для описания свойств кристаллов, квантовых жидкостей (в частности, жидкого гелия), а в дальнейшем при построении ядерных моделей, описании плазмы и т. д.



КВАЗИЭКВИВАЛЕНТНОСТЬ - отношение в классе представлений группы или $C^{*}$-алгебры. А именно, представления $T_{1}$ и $T_{2}$ (группы или $C^{*}$-алгебры) считаются к в а з и э к в и в а ле н т ны м и, если для любого ненулевого подпредставления $T_{1}^{\prime}$ в $T_{1}$ и любого ненулевого подпредставления $T_{2}^{\prime}$ в $T_{2}$ существует ненулевой сплетающий оператор из $T_{1}^{\prime}$ в $T_{2}^{\prime}$



Ю. А. Неретин. KВАНТОВАНИЕ ИЗОТРОПНЫХ ТОРОВ - процедура, сопоставляющая $k$-параметрическому семейству $k$-мерных торов $\Lambda\left(I_{1}, \ldots, I_{k-1}, E\right)$ в $2 n$-мерном фазовом пространстве $M^{2 n}=T^{*} M^{n}$ с комплексным ростком (см. Изотропное подмногообразие с комплексным ростком) (где $\left(I_{1}, \ldots, I_{k-1}\right.$, $E) \in \mathfrak{R} \subset \mathbb{R}^{k}-$ параметры) последовательности асимптотических собственных чисел $E_{v}(h)$ и функщий $\psi_{v}(x, h)$ (спектральную серию) $h$-дифференциального или псевдодифференциального оператора $\hat{L}=L\left(-i h \nabla_{x}, x, h\right)$ в конфигурационном пространстве $M^{n}$, то есть таких $E_{v}$ и $\psi_{\mathrm{v}}$, что



\[

\hat{L} \psi_{v}=E_{v} \psi_{v}+O\left(h^{N}\right), \operatorname{supp} \psi_{v} \rightarrow \pi_{x} \Lambda

\]



при $h \rightarrow 0$. Здесь $N>1, \pi_{x} \Lambda-$ естественная проекция $\Lambda$ в $M^{n}$, и предполагается, что $\left(\Lambda, r^{n}\right)$ инвариантно относительно фазового потока $g_{H}^{+}$, задаваемого гамильтонианом $H=L(p, x, 0)$ (старшей частью символа оператора $\hat{L})$. Числа $E_{v}$, а также соответствующие $E_{v}$ числа $I_{1 v}(h), \ldots, I_{k-1 v}(h)$ из $\mathfrak{R}$ определяются из условий квантования



\[

\frac{1}{2 \pi h} \oint_{\gamma_{j}} \omega^{2}=Q_{j} \bmod 1, \quad j=1, \ldots, k,

\]



где $\gamma_{j}$ - базисные циклы на $\Lambda, \omega^{2}-2$-форма, задающая симплектич. структуру на $M^{2 n}, Q_{j}-$ нек-рые характеристики $\left(\Lambda, r^{n}\right)$. Функции $\psi_{v}$ вычисляются с помощью (комплексного) Маслова канонического оператора (см. [1], [2]).



Если $k=0, \Lambda-$ устойчивая в линейном приближении точка покоя, формулы для $E_{v}, \psi_{v}$ совпадают с формулами осцилляторного приближения.



Если $k=1, \Lambda$ - орбитально устойчивая в линейном приближении замкнутая траектория гамильтоновой системы с гамильтонианом $H$, если $\hat{L} / h^{2} \equiv \Delta-$ оператор Бельтрами Лапласа на римановом многообразии $M^{n}$, то проекция $\Lambda$ на $M^{n}$ - устойчивая в линейном приближении замкнутая геодезическая, $\Psi_{v}$ сосредоточены в ее окрестности (см. [3], [4]).



Если $k=n$ ( $k$ принимает максимальное значение), то $\Lambda-$ (полномерный) лагранжев тор и комплексный росток отсутствует, $O_{j}-$ индекс цикла Маслова $\gamma_{j}$.



Лит.: [1] М а с л о в В. П., Комплексный метод ВКБ в нелинейных уравнениях, М., 1977; [2] Бе ло в В. В., Доб рохото в С. Ю., «Доклады АН СССР», 1988, т. 298, № 5, с. 1037-42; [3] Б а б и ч В. М., Бул д ы р в В. С., Асимптотические методы в задачах дифракции коротких волн, М., 1972; [4] Д о б р ох о то в С. Ю., М а л о в В. П., в кн.: Итоги науки и техники. Современные проблемы математики, т. 23, М., 1983, c. 137-222.



С. Ю. Доброхотов



pq- И qp-KBAНTOBAHИE - квантования алгебры функций на линейном фазовом пространстве $\mathbb{R}^{2 n}$, в которых функциям (символам) вида $p_{1}^{\alpha_{1}} \ldots p_{n}^{\alpha_{n}} q_{1}^{\beta_{1}} \ldots q_{n}^{\beta_{n}}$ ставятся в соответствие операторы $\hat{p}_{1}^{\alpha_{1}} \ldots \hat{p}_{n}^{\alpha_{n}} \hat{q}_{1}^{\beta_{1}} \ldots \hat{q}_{n}^{\beta_{n}}$ в случае $p q$-кван-



\includegraphics[max width=\textwidth, center]{2023_07_08_beb97638694792137dd2g-1}

$\hat{q}_{j}, \hat{p}_{j}$ - операторы координат и импульсов, удовлетворяющие каноническим коммутационным соотношениям:



\[

\left[\hat{q}_{j}, \hat{q}_{k}\right]=\left[\hat{p}_{j}, \hat{p}_{k}\right]=0, \quad\left[\hat{p}_{j}, \hat{q}_{k}\right]=\frac{\hbar}{i} \delta_{i k},

\]



$\hbar-$ постоянная Планка. В $q$-представлении (в пространстве $L^{2}\left(\mathbb{R}_{x}^{n}\right)$ с операторами $\left.\hat{q}_{j}=x_{j}, \hat{p}_{j}=\frac{\hbar}{t} \frac{\partial}{\partial x_{j}}\right)$ общие формулы $p q$-квантования и $q p$-квантования ставят в соответствие функции $f=f(p, q)$ операторы $\hat{f}_{p q}$ и $\hat{f}_{q p}$ вида



\[

\begin{gathered}

\hat{f}_{p q} u(x)=(2 \pi \hbar)^{-n} \int e^{\frac{i}{\hbar} p \cdot(x-y)} f(p, y) u(y) d y \\

\hat{f}_{q p} u(x)=(2 \pi \hbar)^{-n} \int e^{\frac{i}{\hbar} p \cdot(x-y)} f(p, x) u(y) d y .

\end{gathered}

\]



КВАЗИСРЕДНИЕ 243





\end{document}